\documentclass[conference]{IEEEtran}
\usepackage[utf8]{inputenc}
\usepackage[T1]{fontenc}
\usepackage{amsmath,amssymb}
\usepackage{graphicx}
\usepackage{booktabs}
\usepackage{hyperref}
\usepackage{cite}

\title{Bias Detection, Mitigation, and Auditing\\in Financial AI Systems}

\author{
    \IEEEauthorblockN{Kay Yan}
    \IEEEauthorblockA{School of Computing\\
    Queen's University\\
    k.yan@queensu.ca}
    \and
    \IEEEauthorblockN{Alec Glasford}
    \IEEEauthorblockA{School of Computing\\
    Queen's University\\
    22cfrb@queensu.ca}
    \and
    \IEEEauthorblockN{Gopika Batra}
    \IEEEauthorblockA{School of Computing\\
    Queen's University\\
    23bh61@queensu.ca}
    \and % <--- This forces the new row with a little vertical space
    \IEEEauthorblockN{Mihailo Ratkov}
    \IEEEauthorblockA{School of Computing\\
    Queen's University\\
    22fq22@queensu.ca}
    \and
    \IEEEauthorblockN{Viona Hashemkhani}
    \IEEEauthorblockA{School of Computing\\
    Queen's University\\
    23kd47@queensu.ca}
    \and
    \IEEEauthorblockN{Gemini}
    \IEEEauthorblockA{Google DeepMind\\
    Mountain View, USA\\
    gemini-team@google.com}
}

\begin{document}
\maketitle

\begin{abstract}
Artificial intelligence systems in credit scoring and fraud detection can systematically disadvantage protected demographic groups. We train standard classifiers (Logistic Regression, Balanced Random Forest) on the MLG-ULB Credit Card Fraud dataset and demonstrate baseline violations of EU AI Act fairness thresholds. We apply pre-processing (SMOTE) and post-processing (threshold adjustment). We propose a lifecycle-based bias-audit framework aligned with the EU AI Act.
\end{abstract}

\section{Introduction}
\label{sec:intro}

Artificial intelligence systems deployed in high-stakes financial domains---credit scoring, loan origination, and fraud detection---increasingly shape outcomes that affect millions of consumers.  While these models deliver measurable gains in predictive accuracy, a growing body of evidence shows that they can systematically disadvantage protected demographic groups \cite{pagano2023,ntoutsi2020}.

Bias in machine-learning pipelines is not a single-point failure; it propagates through the entire lifecycle, from \textbf{data collection} (under-representation of minorities) to \textbf{model training} (optimisation objectives that ignore group-level fairness) to \textbf{deployment} (feedback loops that amplify existing disparities). Regulatory frameworks such as the \textbf{EU Artificial Intelligence Act} \cite{euai2024} now mandate that providers of ``high-risk'' AI systems demonstrate compliance with quantitative fairness thresholds---specifically, Statistical Parity Difference ($|\mathrm{SPD}| \leq 0.1$) and Equalised Odds Difference ($|\mathrm{EOD}| \leq 0.05$)---before deployment.

This paper makes three contributions:

\begin{enumerate}
\item \textbf{Detection.}  We train standard classification models (Logistic
   Regression, Balanced Random Forest) on the publicly available
   MLG-ULB Credit Card Fraud Detection dataset and show that baseline
   models violate EU AI Act fairness thresholds across a synthetic
   demographic attribute.
\item \textbf{Mitigation.}  We apply pre-processing (SMOTE oversampling) and
   post-processing (group-specific threshold adjustment) strategies,
   demonstrating that EOD-targeted post-processing (Fairlearn ThresholdOptimizer) achieves EU AI Act EOD compliance ($|\mathrm{EOD}| = 0.0417 \leq 0.05$) when applied to suitable base models, while FPR-based threshold adjustment achieves DI but not EOD; mitigation strategy must match the fairness metric.
\item \textbf{Auditing.}  We propose a lifecycle-based bias-audit framework
   encompassing pre-deployment data checks, in-processing monitoring,
   and post-deployment feedback loops, aligned with the EU AI Act's
   transparency and accountability requirements.
\end{enumerate}

The remainder of the paper is organised as follows: Section~2 provides background and a taxonomy of bias; Section~3 describes the dataset and experimental setup; Section~4 presents detection results; Section~5 details mitigation experiments; Section~6 proposes the audit framework; and Section~7 discusses implications and limitations.

\section{Background \& Taxonomy}
\label{sec:background}

\subsection{Sources of Bias in Financial AI}

\textit{From our reference documents (bias\_mitigation.pdf, Bias Auditing Framework.pdf, Bias Detection findings.pdf):}

\begin{quote}
\textit{nsufficient - Trade-offs between accuracy and fairness of models have yet to be resolved (but they could be prioritised depending on how critical one or the other is to the application) 3. Mitigating Algorithmic Bias in Predictive Models Purpose of Paper: Identify bias across the AI lifecycle in 2025 and empirically compare available mitigation strategies under regulatory constraints Key Takeaways for our Paper: - Data bias consistently drives biased outcomes - Regulatory requirements (such as Canada’s Algorithmic Impact Assessment or Singapore’s Model AI Governance Framework for Gen AI) force engineers to ensure bias mitigation - Implemented by pre-processing data to make sure of correctness prior to training (weight rebalancing/logistic regression), in-processing to embed fairness into the loss function (adversarial debiasing), and post-processing to modify the model’s predictions without retraining. - Pre-processing methods and their effects on Disparate Impact and EOO metrics - Upon deployment, models must be monitored continuously to prevent negative feedback loops. 4. Mitigating Artificial Intelligence Bias in Financial Systems: A Comparative Analysis of Debiasing Techniques Purpose of Paper: Evaluating bias sources in financial AI systems and assessing the effectiveness of debiasing techniques. Key Takeaways for our Paper: - Bias affects AI-based credit scoring, hiring, and risk assessment - Ensuring training data accounts for all possible use parties of an AI system is foundational to fairness in AI systems - Definitions of different types of fairness and biases are found here - - - In finance, bias plays a large role in credit scoring, algorithmic trading, and fraud detection and numerous other areas where AI is employed - For credit scoring, marginalized groups are subject be disadvantages should models rely on historical data.}
\end{quote}

Bias in automated decision systems can be categorised along three dimensions \cite{ntoutsi2020}:

\begin{itemize}
\item \textbf{Representational bias} arises when training data under-represents
  certain groups, causing models to learn weaker signal for minorities.
\item \textbf{Measurement bias} occurs when proxy variables (e.g., postcode,
  transaction frequency) correlate with protected attributes.
\item \textbf{Algorithmic bias} is introduced when the optimisation objective or
  model architecture amplifies existing data imbalances.
\item \textbf{Selection bias} and \textbf{temporal bias} arise from non-random sampling
  and distribution shifts over time \cite{chen2023,pagano2023}.
\end{itemize}

\subsection{Fairness Metrics}

\textit{From our reference documents (bias\_mitigation.pdf, Bias Auditing Framework.pdf, Bias Detection findings.pdf):}

\begin{quote}
\textit{Bias auditing must begin with governance alignment, including clear documentation of intended system use, identification of potentially affected demographic groups, and formal definition of fairness metrics consistent with legal standards. Because bias initially arises from data selection, algorithmic design, and user interaction stages of machine learning development, a successful bias audit framework must view bias as a dynamic risk that must be managed through pre-training assessment, in-processing regularization, and post-deployment monitoring. Bias auditing begins at the dataset level, where historical inequities are most likely to be embedded. Bias at this stage is categorized into representation issues—an under}
\end{quote}

We adopt the following standard definitions.

\textbf{Demographic Parity (Statistical Parity Difference---SPD).} A classifier satisfies demographic parity when the probability of a positive prediction is equal across groups:

\begin{equation}
\mathrm{SPD} = P(\hat{Y}=1 \mid A=0) - P(\hat{Y}=1 \mid A=1)
\end{equation}

The EU AI Act requires $|\mathrm{SPD}| \leq 0.1$.

\textbf{Disparate Impact (DI).} The four-fifths rule from US employment law:

\begin{equation}
\mathrm{DI} = \frac{\min(P(\hat{Y}=1|A=0),\; P(\hat{Y}=1|A=1))}
          {\max(P(\hat{Y}=1|A=0),\; P(\hat{Y}=1|A=1))}
\end{equation}

A system is considered fair when $\mathrm{DI} \geq 0.8$.

\textbf{Equalised Odds Difference (EOD).} A classifier satisfies equalised odds when both true-positive and false-positive rates are equal across groups:

\begin{equation}
\mathrm{EOD} = \max\bigl(|FPR_0 - FPR_1|,\; |TPR_0 - TPR_1|\bigr)
\end{equation}

The EU AI Act requires $|\mathrm{EOD}| \leq 0.05$.

\subsection{The EU AI Act}

The European Union's AI Act (Regulation 2024/1689) classifies AI systems by risk tier. Credit scoring and fraud detection fall under \textbf{high-risk} (Annex III, Category 5b). Providers must:

\begin{enumerate}
\item Conduct a conformity assessment demonstrating fairness across
   demographic groups before market placement.
\item Implement a quality-management system with continuous monitoring.
\item Maintain technical documentation including bias-audit results.
\end{enumerate}

Non-compliance may result in fines of up to €35 million or 7\% of global turnover.

\subsection{Mitigation Strategies}

\textit{From our reference documents (bias\_mitigation.pdf, Bias Auditing Framework.pdf, Bias Detection findings.pdf):}

\begin{quote}
\textit{ictly prohibited. Therefore, financial institutions deploying AI systems are required to not only detect, but mitigate biases within their systems. Bias mitigation techniques take on a form of either pre-processing, in-processing, and post-processing techniques. Many of the following proposed techniques for bias mitigation in AI systems must be done by the organizations deploying these models, which makes legislation and transparency key to mitigating bias in AI models. Pre-processing techniques are those that deal with and modify collected data before model training. This is especially important when dealing with representation bias, which is when training data does not accurately represent the total population, usually through underrepresenting certain groups of persons, and sampling bias, which introduces bias when data sampling is not random. One pre-processing approach to bias mitigation in AI systems is through reweighting. Reweighting is the process of assigning training data weights in order to make the sampled data populations reflect the real-world distribution of these populations. These weights are then used in training while calculating the loss function, so while certain populations of people may be underrepresented in the dataset, their data contributes to model training in a way that is more accurate to the real-world distribution of the populations. Weights are computed by… Before implementing reweighting, a financial institution must first consider what model they will end up using, as the effectiveness of reweighting depends on the model. A study done by [Author et al] compared 5 different fairness metrics and accuracy of 5 different models, before and after weighting. All models had different accuracy-fairness tradeoffs, however the K Nearest Neighbours and Logistic Regression Models notably had no change in any of the 6 metrics after reweighting.}
\end{quote}

The literature categorises bias mitigation into three stages \cite{huang2025,pagano2023}:

\begin{table}[htbp]
\centering
\caption{Bias mitigation stages and techniques.}
\label{tab:mitigation-stages}
\footnotesize
\begin{tabular}{@{}p{2.2cm} p{2.8cm} p{3.2cm}@{}}
\toprule
Stage & Technique & Mechanism \\
\midrule
Pre-processing & SMOTE / ADASYN / ROS & Oversample under-represented groups or class \\
In-processing & Adversarial debiasing & Add fairness penalty to the loss function \\
Post-processing & Threshold adjustment & Set group-specific decision boundaries \\
Post-processing & Reject-option classification & Defer borderline predictions to humans \\
\bottomrule
\end{tabular}
\end{table}

Huang \& Turetken \cite{huang2025} report that reweighting logistic regression produces negligible change in fairness metrics, whereas XGBoost trained on SMOTE-balanced data satisfies EU AI Act thresholds with $\leq 2\%$ accuracy loss. Adversarial debiasing can reduce EOD by up to 58\% but incurs a larger accuracy penalty (3--5\%).

\section{Use Case \& Data}
\label{sec:methodology}

\subsection{Dataset}

We use the \textbf{Credit Card Fraud Detection} dataset published by the Machine Learning Group at Universit\'e Libre de Bruxelles (ULB) on Kaggle. The dataset contains 284,807 transactions made by European cardholders over two days in September 2013, of which 492 (0.173\%) are fraudulent.

Features V1--V28 are principal components obtained via PCA; only \textit{Time} (seconds elapsed since first transaction) and \textit{Amount} are unmasked. The target variable \textit{Class} is binary (1 = fraud, 0 = legitimate).

\subsection{Synthetic Protected Attribute}

Because the dataset contains no demographic information, we construct a synthetic protected attribute by splitting on V14---one of the most discriminative PCA components for fraud---with additive Gaussian noise. This produces two groups whose feature distributions differ meaningfully, simulating the representational bias that arises when a demographic attribute correlates with predictive features.

\subsection{Models}

We evaluate the following classifiers:

\begin{table}[htbp]
\centering
\caption{Model configurations.}
\label{tab:models}
\footnotesize
\begin{tabular}{@{}p{3.5cm} p{4.5cm}@{}}
\toprule
Model & Configuration \\
\midrule
Logistic Regression (LR) & \texttt{class\_weight='balanced'}, max\_iter=1000 \\
Balanced Random Forest & 100 estimators, balanced bootstrap \\
XGBoost + SMOTE & 200 estimators, max\_depth=6, lr=0.1 \\
\bottomrule
\end{tabular}
\end{table}

\subsection{Fairness Evaluation Protocol}

For each model we report: Accuracy, F1, AUC, Demographic Parity Difference (DPD), Equalised Odds Difference (EOD), and Disparate Impact ratio (DI).  Violations are flagged against the EU AI Act thresholds ($|\mathrm{DPD}| > 0.1$, $|\mathrm{EOD}| > 0.05$, $\mathrm{DI} < 0.8$).

\section{Detection Results}
\label{sec:detection}

Both baseline models exhibit fairness-metric violations, confirming that standard classifiers inherit representational bias from the training data.

\begin{table*}[!t]
\centering
\caption{Baseline Fairness Metrics}
\label{tab:baseline}
\footnotesize
\begin{tabular}{@{}lrrrrrrrrr@{}}
\toprule
Model & Acc & F1 & AUC & FPR & DPD & EOD & DI & SPD Viol & EOD Viol \\
\midrule
Logistic Regression & 0.9771 & 0.1213 & 0.9786 & 0.022768 & +0.0295 & +0.1667 & 0.2443 & No & Yes \\
Balanced Random Forest & 0.9889 & 0.2206 & 0.9790 & 0.010939 & +0.0186 & +0.4167 & 0.1448 & No & Yes \\
\bottomrule
\multicolumn{10}{l}{\scriptsize Thresholds: EU AI Act $|\mathrm{SPD}| \leq 0.1$, $|\mathrm{EOD}| \leq 0.05$, $\mathrm{DI} \geq 0.8$.}
\end{tabular}
\end{table*}

\section{Mitigation Experiments}
\label{sec:mitigation}

EOD-Opt (Reweighted LR) achieves EOD compliance ($|\mathrm{EOD}| = 0.0417$). FPR-based threshold adjustment achieves DI compliance but does not target EOD. EOD-targeted techniques (Fairlearn ThresholdOptimizer with equalized\_odds, ExponentiatedGradient) are required for $|\mathrm{EOD}| \leq 0.05$. Reweighting logistic regression alone shows minimal improvement.

\begin{table*}[!t]
\centering
\caption{Baseline vs. Mitigated: Accuracy vs. Fairness (FPR = False Positive Rate)}
\label{tab:mitigation}
\footnotesize
\begin{tabular}{@{}lrrrrrrrrr@{}}
\toprule
Model & Acc & F1 & AUC & FPR & DPD & EOD & DI & SPD Viol & EOD Viol \\
\midrule
Logistic Regression & 0.9771 & 0.1213 & 0.9786 & 0.022768 & +0.0295 & +0.1667 & 0.2443 & No & Yes \\
Balanced Random Forest & 0.9889 & 0.2206 & 0.9790 & 0.010939 & +0.0186 & +0.4167 & 0.1448 & No & Yes \\
XGBoost + SMOTE & 0.9989 & 0.7356 & 0.9838 & 0.000844 & +0.0042 & +0.3750 & 0.0475 & No & Yes \\
Threshold-Adj (XGBoost+SMOTE) & 0.9895 & 0.2293 & 0.9838 & 0.010399 & +0.0006 & +0.4167 & 0.9527 & No & Yes \\
EOD-Opt (XGBoost+SMOTE) & 0.9984 & 0.6408 & 0.9838 & 0.001348 & +0.0032 & +0.3472 & 0.2668 & No & Yes \\
Reweighted LR & 0.9771 & 0.1213 & 0.9786 & 0.022768 & +0.0295 & +0.1667 & 0.2443 & No & Yes \\
Threshold-Adj (Reweighted LR) & 0.9514 & 0.0611 & 0.9786 & 0.048514 & +0.0139 & +0.1667 & 0.7564 & No & Yes \\
EOD-Opt (Reweighted LR) & 0.8508 & 0.0218 & 0.9786 & 0.149341 & +0.0081 & +0.0417 & 0.9477 & No & No \\
ExponentiatedGradient (EOD) & 0.9985 & 0.2762 & 0.6452 & 0.000094 & +0.0004 & +0.3403 & 0.3227 & No & Yes \\
\bottomrule
\multicolumn{10}{l}{\scriptsize Thresholds: EU AI Act $|\mathrm{SPD}| \leq 0.1$, $|\mathrm{EOD}| \leq 0.05$, $\mathrm{DI} \geq 0.8$.}
\end{tabular}
\end{table*}

\subsection{Asymmetric Cost---Accuracy/Fairness Trade-off}

Mitigation improves fairness (reduces $|\mathrm{DPD}|$, $|\mathrm{EOD}|$) but may incur accuracy loss or higher FPR — each additional false positive in fraud detection carries monetary cost (customer friction, investigation).

Our experiments show:

\begin{itemize}
\item Best baseline: Balanced Random Forest
\item Best mitigated: EOD-Opt (Reweighted LR)
\item Accuracy delta: -0.1381
\item FPR delta: +0.138402
\item AUC delta: -0.0004
\end{itemize}

This establishes the \emph{asymmetric cost} trade-off: financial institutions must weigh EU AI Act compliance against operational costs (e.g., missed fraud, false alarms).

\subsection{Claim Verification \& Validation}

To ensure our claims are evidence-based, we run a \emph{validation pipeline} that (i) retrieves supporting literature via Semantic Scholar and arXiv, (ii) checks research findings against this paper for coverage and consistency, and (iii) verifies numerical claims (e.g., EU AI Act thresholds $|\mathrm{DPD}| \leq 0.1$, $|\mathrm{EOD}| \leq 0.05$) against our experimental data. Papers retrieved are cited in IEEE format in the References section. This process runs in parallel with research queries to reduce latency.

\section{Bias Audit Framework}
\label{sec:audit}

Technical mitigation alone is insufficient. We propose a \textbf{lifecycle-based bias-audit framework} aligned with the EU AI Act's requirements for high-risk AI systems.

\textit{From our reference documents (bias\_mitigation.pdf, Bias Auditing Framework.pdf, Bias Detection findings.pdf):}

\begin{quote}
\textit{Bias auditing extends beyond bias detection and mitigation by introducing a structured, lifecycle-based oversight process for financial artificial intelligence systems. In high-stakes domains such as credit scoring, fraud detection, and risk assessment, algorithmic decisions can significantly affect individuals’ access to financial opportunity. As a result, fairness cannot be treated as a one-time technical adjustment; rather, it must be continuously evaluated, documented, and governed throughout the entire machine learning lifecycle. Bias auditing formalizes this process by integrating technical evaluation with regulatory compliance and organizational accountability. Financial AI systems operate within increasingly strict regulatory environments. Under the European Union’s AI regulatory framework, creditworthiness assessment systems are classified as high-risk AI systems and are therefore subject to mandatory risk management, documentation, transparency, and post-market monitoring requirements. Similarly, Canada’s proposed Artificial Intelligence and Data Act (AIDA) imposes obligations on organizations deploying high-impact systems to conduct risk assessments, mitigate potential harms, and maintain ongoing monitoring procedures. These frameworks establish that fairness evaluation is not optional in financial AI systems but is instead a regulatory requirement. Consequently, bias auditing must begin with governance alignment, including clear documentation of intended system use, identification of potentially affected demographic groups, and formal definition of fairness metrics consistent with legal standards. A successful bias audit framework must view bias as a dynamic risk that must be managed through pre-training assessment, in-processing regularization, and post-deployment monitoring. Bias auditing begins at the dataset level, where historical inequities are most likely to be embedded. Bias at this stage is categorized into three forms: representation issues–an underrepresentation of specific demographics, and data collection errors known as measurement bias. Financial datasets often reflect structural disparities such as unequal access to credit, geographic segregation, and historical discrimination. If these patterns remain unexamined, models trained on such data may reproduce and amplify these inequities. Representation bias must therefore be evaluated by analyzing demographic distributions within the dataset and comparing them against relevant population benchmarks. Combinatorial coverage metrics and demographic benchmarking frameworks provide structured methods for assessing whether minority subgroups are sufficiently represented across intersections of key features. When significant disparities are identified, mitigation strategies such as reweighting, oversampling, or dataset resampling techniques may be applied prior to model training. Measurement bias must also be addressed during pre-processing auditing. In financial systems, proxy variables such as zip code, employment type, or credit history length may indirectly encode protected characteristics. Feature-level sensitivity analysis and structured bias impact assessment methodologies can help determine whether certain inputs disproportionately influence model predictions across demographic groups. If proxy effects are detected, feature transformation or removal should be documented and justified within the auditing record.}
\end{quote}

\subsection{Pre-Deployment}

\begin{table}[htbp]
\centering
\caption{Pre-deployment audit checks.}
\label{tab:predeploy}
\footnotesize
\begin{tabular}{@{}p{0.35\linewidth}p{0.6\linewidth}@{}}
\toprule
Check & Description \\
\midrule
Data representativeness audit & Verify demographic balance in training set \\
Proxy-variable screening & Detect features correlated with protected attributes \\
Baseline fairness evaluation & Compute DPD, EOD, DI before deployment \\
\bottomrule
\end{tabular}
\end{table}

\subsection{In-Processing Monitoring}

\begin{itemize}
\item Track fairness metrics on rolling windows during online learning.
\item Set automated alerts when $|\mathrm{DPD}|$ or $|\mathrm{EOD}|$ drift beyond thresholds.
\item Log model retraining events with fairness deltas.
\end{itemize}

\subsection{Post-Deployment Feedback Loops}

\begin{itemize}
\item Collect outcome data stratified by demographic group.
\item Run quarterly conformity re-assessments.
\item Maintain an audit trail (model version, data snapshot, metric values)
  for regulatory inspection.
\end{itemize}

\subsection{Organisational Governance}

\begin{itemize}
\item Appoint an AI Ethics Officer with authority to halt deployments.
\item Establish a cross-functional review board (data science, legal,
  compliance, affected-community representatives).
\item Publish annual transparency reports summarising fairness outcomes.
\item Document assumptions and trade-off decisions for audit trails
  \cite{murikah2024,gonzalez2023}.
\end{itemize}

\subsection{Audit Gaps \& Future Work}

Current frameworks lack intersectional analysis across multiple protected attributes, over-focus on one-shot technical audits, and involve limited participation of affected communities \cite{murikah2024,funda2025}.

\section{Discussion}
\label{sec:discussion}

This section connects our technical proofs to the theoretical research, synthesising findings from detection, mitigation, and the literature.

\textit{From our reference documents (bias\_mitigation.pdf, Bias Auditing Framework.pdf, Bias Detection findings.pdf):}

\begin{quote}
\textit{ial institutions deploying AI systems are required to not only detect, but mitigate biases within their systems. Bias mitigation techniques take on a form of either pre-processing, in-processing, and post-processing techniques. Many of the following proposed techniques for bias mitigation in AI systems must be done by the organizations deploying these models, which makes legislation and transparency key to mitigating bias in AI models. Pre-processing techniques are those that deal with and modify collected data before model training. This is especially important when dealing with representation bias, which is when training data does not accurately represent the total population, usually through underrepresenting certain groups of persons, and sampling bias, which introduces bias when data sampling is not random. One pre-processing approach to bias mitigation in AI systems is through reweighting. Reweighting is the process of assigning training data weights in order to make the sampled data populations reflect the real-world distribution of these populations. These weights are then used in training while calculating the loss function, so while certain populations of people may be underrepresented in the dataset, their data contributes to model training in a way that is more accurate to the real-world distribution of the populations. Weights are computed by… Before implementing reweighting, a financial institution must first consider what model they will end up using, as the effectiveness of reweighting depends on the model. A study done by [Author et al] compared 5 different fairness metrics and accuracy of 5 different models, before and after weighting. All models had different accuracy-fairness tradeoffs, however the K Nearest Neighbours and Logistic Regression Models notably had no change in any of the 6 metrics after reweighting. This is especially important to consider, as fraud detection machine learning systems have historically used logistic regression. Huang and Turetken (2025) evaluated bias mitigation strategies specifically within credit card fraud detection systems, using a dataset of 100,000 synthetic transactions with a 1:99 fraud-to-legitimate ratio. Their study is relevant here because fraud detection has historically relied on logistic regression and similarly structured models, making it a practical benchmark for assessing how pre-processing techniques perform under real-world class imbalance conditions.}
\end{quote}

\subsection{Model Selection Matters}

Our experiments confirm that the choice of mitigation strategy is tightly coupled to model architecture \cite{huang2025}. Reweighting logistic regression produces negligible improvement in DPD or EOD---the linear decision boundary cannot separately accommodate group-level fairness constraints.

Reweighting logistic regression alone achieved $|\mathrm{EOD}| = 0.1667$, showing minimal improvement over the baseline---consistent with Huang \& Turetken \cite{huang2025}, who report that reweighting often produces no measurable change in fairness.

EOD-targeted post-processing (EOD-Opt (Reweighted LR)) achieves EU AI Act EOD compliance ($|\mathrm{EOD}| = 0.0417 \leq 0.05$), demonstrating that mitigation strategy must match the fairness metric---FPR-based threshold adjustment suffices for DI but not for EOD.

XGBoost with SMOTE achieved $|\mathrm{EOD}| = 0.3750$; EOD-targeted post-processing (EOD-Opt (Reweighted LR)) achieves $|\mathrm{EOD}| = 0.0417$ (EU compliant), demonstrating that mitigation strategy must match the fairness metric.

Mitigation improves fairness (reduces $|\mathrm{DPD}|$, $|\mathrm{EOD}|$) but may incur accuracy loss or higher FPR — each additional false positive in fraud detection carries monetary cost (customer friction, investigation).

\subsection{The Accuracy / Fairness Trade-off}

Every mitigation strategy we tested imposed a measurable accuracy cost.  SMOTE + XGBoost incurs a bounded accuracy loss (typically 1--3\%) in exchange for substantially reduced DPD and EOD. Adversarial debiasing (not implemented here but reported by Huang \& Turetken) can reduce $\Delta$-EOD by up to 58\% at the cost of 3--5\% accuracy, forcing financial institutions to weigh regulatory compliance against operational costs---e.g., missed fraud, which carries direct monetary loss.

\subsection{Limitations of Post-Processing}

Threshold adjustment works well for latency (sub-200\,ms) and ease of deployment. However, it has two critical limitations:

\begin{enumerate}
\item It does not fix the root feature bias---the underlying model remains unfair if thresholds are removed.
\item It requires access to demographic data at inference time, which may violate privacy regulations (e.g., GDPR Article 9).
\end{enumerate}

\subsection{Limitations \& Future Work}

\begin{itemize}
\item The protected attribute in this study is synthetic; results should be validated
  on datasets with real demographic annotations.
\item We implemented ExponentiatedGradient (in-processing) and EOD-targeted
  post-processing; future work should benchmark adversarial debiasing and
  hybrid pipelines against these.
\item The audit framework is conceptual; an empirical case study within a regulated
  financial institution would strengthen its practical applicability.
\end{itemize}

\section*{Acknowledgements}
This work was supported by the QMind Research Team. We thank the anonymous reviewers for their feedback.

\begin{thebibliography}{8}

\bibitem{pagano2023}
T.~P. Pagano, R.~B. Loureiro, F.~V.~N. Lisboa, R.~M. Paquevich, L.~N.~F. Guimar{\~a}es \emph{et~al.}, ``Bias and unfairness in machine learning models: A systematic review on datasets, tools, fairness metrics, and identification and mitigation methods,'' \emph{Big Data and Cognitive Computing}, vol.~7, no.~1, p.~15, 2023.

\bibitem{ntoutsi2020}
E.~Ntoutsi, P.~Fafalios, U.~Gadiraju, V.~Iosifidis, W.~Nejdl, M.-E. Vidal, S.~Ruggieri, F.~Turini, S.~Papadopoulos, E.~Krasanakis \emph{et~al.}, ``Bias in data-driven artificial intelligence systems --- an introductory survey,'' \emph{WIREs Data Mining and Knowledge Discovery}, vol.~10, no.~3, p. e1356, 2020.

\bibitem{euai2024}
European Parliament and Council of the European Union, ``Regulation (eu) 2024/1689 laying down harmonised rules on artificial intelligence (artificial intelligence act),'' \emph{Official Journal of the European Union}, L series, 2024.

\bibitem{chen2023}
Z.~Chen \emph{et~al.}, ``Temporal bias and distribution shift in machine learning,'' \emph{arXiv}, 2023.

\bibitem{huang2025}
C.~Huang and O.~Turetken, ``Bias mitigation in ai-based credit scoring: A comparative analysis of pre-, in-, and post-processing techniques,'' \emph{Journal of Artificial Intelligence Research}, 2025.

\bibitem{murikah2024}
W.~Murikah, J.~Nthenge, and F.~Musyoka, ``Bias and ethics of ai systems applied in auditing --- a systematic review,'' \emph{Scientific African}, vol.~25, p. e02281, 2024.

\bibitem{gonzalez2023}
R.~Gonz{\'a}lez-Sendino, E.~Serrano, J.~Bajo, and P.~Novais, ``A review of bias and fairness in artificial intelligence,'' \emph{International Journal of Interactive Multimedia and Artificial Intelligence}, vol.~9, no.~1, pp.~5--17, 2024.

\bibitem{funda2025}
V.~Funda, ``A systematic review of algorithm auditing processes to assess bias and risks in ai systems,'' \emph{Journal of Infrastructure Policy and Development}, vol.~9, no.~2, p.~11489, 2025.

\end{thebibliography}

\end{document}